\section{Endpoint Interfaces and Certificate Reuse}\label{ta:endpoint-interface}
Let $K$ and $K'$ satisfy Theorem~\ref{thm:endpoint-interface}. Their local data define the same finite ambient state space of version-tagged component tuples, monitor memories and pending-command sets. Let $I,I'$ be their entry sets, and let $T,T'$ be their target projections. A target projection retains all component and new-monitor coordinates and removes only the endpoint-controller coordinate. All targets retain their eligibility premises, including reachability and endpoint safety, deadlock freedom and preservation of uncontrollable events.

\paragraph{Identical interfaces.}
Assume $I=I'$ and $T=T'$. On every ambient state, enabled ordinary events and their complete outcome buckets are equal because the component models and controllability agree. Transfer guards and outcomes agree because the transfer relations, pending commands and precedence agree. Boundary guards and memory updates agree because monitors and initializers agree. UC quiescence and safety therefore agree as well. Induction on path length from $I$ gives equal reachable state sets and label/outcome transitions. No rule inspects either endpoint-controller state after entry. Old states collapsed by $\Init$ may have different controller coordinates, but such coordinates impose no source restriction once old control is released.

A reached state is a goal precisely when it is safe, UC-quiescent, has no pending command, and its component/new-monitor projection belongs to $T$. Hence goals agree. The identity map preserves the policy successor relation, safety, nonblocking, all maximal paths and every strict rank inequality. It also preserves the losing-certificate obligations: goals are absent from $L$, a losing entry remains in $L$, and at each safe member the same UC escape witness or all-control counter-outcomes remain in $L$. Thus WIN and LOSS certificates transfer, including their universal quantification over all enabled outcomes.

For every reached goal $q$, membership of its projection in $T'$ supplies at least one full target $z'\in Z'_{\mathrm{load}}$. Choose one and define $\kappa'(q)=z'$. The selected endpoint controller is eligible and the handover copies precisely its component, monitor and controller coordinates. Proposition~\ref{prop:endpoint-continuation} yields an isomorphism to that endpoint rooted at $z'$. This preserves that endpoint's complete continuation language; it does not claim equivalence between the two endpoint controllers' continuation languages. Nondeterministic local outcomes are unchanged and remain adversarial. Several targets with one projection cause no difficulty: selecting a compatible controller state is part of handover, not selection among an event's nondeterministic outcomes.

\paragraph{A Cell replacement with different continuation behavior.}
Delete only $c_0\xrightarrow{j}c_0$ from the new controller in Table~\ref{tab:cell-finite-main}. Its alphabet still contains $j$ through $c_1\xrightarrow{j}c_0$. The reachable endpoint tuples remain exactly $z_A=(h,e,c_0,c)$, $y_B=(e,h,c_1,d)$ and $z_B=(e,h,c_0,c)$: the cycle $z_A\xrightarrow{m_B}y_B\xrightarrow{j}z_B\xrightarrow{m_A}z_A$ remains, and the deleted edge removes only the inspection self-loop at $z_B$. This closed loop remains safe and deadlock-free; all events are controllable. Keeping $Z_{\mathrm{load}}=\{z_A,z_B\}$ preserves its target projection, and the old endpoint is unchanged. Thus Figure~\ref{fig:cell-story}'s whole update certificate and its handover choices transfer. The endpoint languages differ: from $z_B$, the original admits $j^\omega$, whereas the replacement's first event must be $m_A$. Reuse therefore does not require continuation-language equivalence between endpoints. This finite construction adds no solver measurement.

\paragraph{Smaller entry interface and larger target interface.}
Assume instead $I'\subseteq I$ and $T\subseteq T'$, and let $(\sigma,\rank,\kappa)$ win for $K$. Starting from $I'$, use $\sigma$ until the first goal of $K'$ and stop there. Every resulting nonterminal prefix is also a prefix of the original policy from $I$, because all local transitions agree. Safety and nonblocking are inherited. An infinite resulting path would be an infinite original-policy path, contradicting its finite completion. Every original goal has its projection in $T'$ and is also a goal when reached from $I'$, so a maximal resulting path cannot pass an old goal without stopping. A new goal may be reached earlier, which only truncates a path.

Restrict $\rank$ to this truncated closure and assign zero at its goals. Edges between remaining non-goals keep their original strict inequalities. An edge entering a goal has a positive-rank non-goal source and zero-rank target, so also decreases. Choose a compatible target in $Z'_{\mathrm{load}}$ at each such goal. This supplies a complete WIN certificate for $K'$. The converse is not claimed: fewer entries or additional targets can make a formerly losing contract win. Nor does the theorem preserve the particular policy returned by a new solver run, which may use a different endpoint-dependent exploration order.
